\documentclass[11pt]{article}
\usepackage[hmargin=2.5cm,vmargin=2.5cm]{geometry}
\usepackage{amsmath,amssymb}
\usepackage[colorlinks=true,linkcolor=blue,urlcolor=blue]{hyperref}

\title{The Intermediate Firm's Cost-Minimization Problem:\\ Factor Demands and the Markup Wedge}
\author{Spending-efficiency model --- internal note}
\date{\today}

\begin{document}
\maketitle

\noindent
This note derives the factor-demand conditions of Appendix~A in the working paper
(\texttt{../../draftPaper.tex}) from the intermediate firm's cost-minimization
problem, and shows where the fixed gross markup $\mu^{p}$ (model parameter
\texttt{markupss}) enters. It also records a fixed-cost technicality that is inert
at the current calibration.

\section{The problem}

Firm $i$ takes the factor prices $r_t^{k}$ and $w_t$ and the externalities
$K_{t-1}^{GI}$ (public infrastructure) and $A_{t-1}$ (adopted technology) as given,
and chooses private capital and effective labor to produce a required output level
$\bar{Y}_t(i)$ at minimum factor cost:
\begin{equation}\label{eq:problem}
    \min_{K_{t-1}(i),\,N_t(i)} \; r_t^{k}\,K_{t-1}(i) + w_t\,N_t(i)
\end{equation}
subject to the production technology
\begin{equation}\label{eq:constraint}
    A_{t-1}^{\vartheta-1}\left(K_{t-1}^{GI}\right)^{\alpha_G}
    K_{t-1}(i)^{\alpha}\,N_t(i)^{1-\alpha} - \Theta \;\ge\; \bar{Y}_t(i).
\end{equation}
Stating the required output level is what makes the problem well-posed: without it,
minimizing cost subject to the technology alone has the degenerate zero-input
solution.

\section{Lagrangian and first-order conditions}

With the multiplier $\xi_t$ on the output constraint,
\begin{equation}
    \mathcal{L} = r_t^{k}\,K_{t-1}(i) + w_t\,N_t(i)
    + \xi_t\left[\bar{Y}_t(i)
    - A_{t-1}^{\vartheta-1}\left(K_{t-1}^{GI}\right)^{\alpha_G}
      K_{t-1}(i)^{\alpha}\,N_t(i)^{1-\alpha} + \Theta\right].
\end{equation}
The first-order conditions are
\begin{align}
    [\,K_{t-1}(i)\,]:&\quad r_t^{k} = \xi_t\,\alpha\,\frac{Y_t+\Theta}{K_{t-1}(i)},
    \label{eq:focK}\\
    [\,N_t(i)\,]:&\quad w_t = \xi_t\,(1-\alpha)\,\frac{Y_t+\Theta}{N_t(i)},
    \label{eq:focN}
\end{align}
where $Y_t$ denotes output net of the fixed cost, so $Y_t+\Theta$ is the gross
Cobb--Douglas term whose marginal products appear in the conditions. Dividing
\eqref{eq:focK} by \eqref{eq:focN} gives the factor-mix condition
\begin{equation}\label{eq:mix}
    \frac{K_{t-1}(i)}{N_t(i)} = \frac{\alpha}{1-\alpha}\,\frac{w_t}{r_t^{k}},
\end{equation}
which is common across firms, since all face the same factor prices and technology
(model equation \texttt{Kp(-1)/N = alpha/(1-alpha)*w/rk},
\texttt{models/model\_block.modpart}).

\section{The multiplier and the markup wedge}

By the envelope theorem, $\xi_t$ is the shadow cost of one more unit of output: the
firm's \emph{real marginal cost}. With constant returns to the two chosen inputs, it
also equals unit factor cost.

The firm sells its output to the Calvo retailers at the fixed gross markup $\mu^{p}$
over this cost. The resulting relative price of intermediate output is $mc_t$ (the
retailers' marginal cost), so
\begin{equation}\label{eq:wedge}
    mc_t = \mu^{p}\,\xi_t
    \qquad\Longleftrightarrow\qquad
    \xi_t = \frac{mc_t}{\mu^{p}}.
\end{equation}
Substituting \eqref{eq:wedge} into \eqref{eq:focK}--\eqref{eq:focN} (with
$\Theta=0$; see Section~\ref{sec:theta}) gives the paper's factor demands,
\begin{equation}\label{eq:factordemand}
    \boxed{\;r_t^{k} = \alpha\,\frac{mc_t}{\mu^{p}}\,\frac{Y_t}{K_{t-1}}, \qquad
    w_t = (1-\alpha)\,\frac{mc_t}{\mu^{p}}\,\frac{Y_t}{N_t}\;}
\end{equation}
matching the model's labor-demand equation
\texttt{(1-alpha)*mc*y/N = markupss*w} and the steady state
\texttt{w=(1-alpha)*mc*y/N/markupss}
(\texttt{models/Steady\_states\_solution.m}).

Factor payments therefore absorb only $mc_t/\mu^{p}$ of revenue per unit of output.
The residual,
\begin{equation}
    \frac{\mu^{p}-1}{\mu^{p}}\,mc_t\,Y_t,
\end{equation}
is the profit flow that accrues to the owners of adopted technologies --- the
dividend term in the value function $\mathcal{V}_t$
(model equation \texttt{V = (markupss-1)/(markupss)*mc*y + ...}). Without the wedge
in \eqref{eq:factordemand}, this profit flow would have no source: competitive
factor pricing at $mc_t$ would exhaust the entire revenue.

Note that $\mu^{p}$ (\texttt{markupss} $=1.18$) is a separate calibrated parameter,
\emph{not} the retail Calvo markup $\epsilon/(\epsilon-1)=10/9\approx 1.11$. The two
markups belong to different layers: $\mu^{p}$ to the intermediate producers, whose
profits reward technology adoption; $\epsilon$ to the retailers, whose staggered
pricing generates the New Keynesian block.

\section{The fixed-cost technicality}\label{sec:theta}

The first-order conditions \eqref{eq:focK}--\eqref{eq:focN} involve the gross
Cobb--Douglas term $Y_t+\Theta$, while the model's factor-demand equations use
$Y_t$ (net of the fixed cost). The two coincide because the fixed cost is switched
off at the calibration: \texttt{Bigtheta=0}, \texttt{Bigtheta\_y=0}
(\texttt{models/parameters\_common.macro}). If $\Theta$ were ever calibrated
positive, the model's factor demands would need \texttt{y+Bigtheta} in place of
\texttt{y} to remain the exact first-order conditions.

\end{document}
